% Nouveau EGA V (cote 25): archive cote 25, site page 94 (PDF page 95).
% typescript. Draft compared: none.
% AI-transcribed by Claude Opus 5.5 (claude-opus-5-5, adaptive thinking, effort high), Message Batches API, batch ega25b (msgbatch_01A7n5YuXubUPNK1zsepi8RG), 2026-11-13; not yet checked by eye.
\scanpage{25/94}
\mspage{88}
semble qu'elle doive faire un axiome primitif de la situation (dans la
situation du n°2.2. , elle \unsure{résultait} d'un argument global assez \unsure{spécifil|spécifiq} ):
\uline{pour toute partie fermée Zrréductible} Z \uline{de} P d\uline{e dimension} m, $\underline{Z}_{\hand{\eta}} \neq \emptyset$ \xout{xxx}
\xout{pour tout La précision sur la dimension sera alors conséquence facile}
\xout{de l'hypothèse suivante, de nature différentielle: P est lisse sur S,}
\xout{de dimension $\underline{P}$ $P \times_S T$ une immersion, et pour tout} Reprenant alors la
partie fermée Z de X,\xout{de di} telle que dim f(Z) $\hand{\geqslant}$ m, on voit que $\underline{Z}_{\hand{\eta}} \rightarrow G$ est
dominant, et par suite $\underline{Z}_{\hand{\eta}}$ est de dimension dim $\underline{Z}$ - dim G = dim Z-m.

Ces propriétés permettent de développer dans le contexte présent
les résultats correspondant à 2.1. à 2.11. Il faut une condition sur (D)
assurant la validité de 2.12. i.e. que si X est lisse, alors $\underline{X}$ l'est
\unsure{éggalement|également}, en supposant f:X $\longrightarrow$ P non ramifié. N$_{\mathrm{o}}$us supposerons
\unsure{maàntenant|maintenant} P \uline{lisse sur} k,\inserted{$\underline{P} \rightarrow P \times_k G$ quasi-fini,} et la condition suivante vérifiée\inserted[typed]{(où nous supposons k alg clos)}: pour
tout $x \in P(k)$, et tout sous-espace vectoriel V \inserted[typed]{de dimension $n \geqslant m$} de l'espace tangent
$T_x$(P) de P en x, considérons l'ensemble\inserted{E(x,V)} \xout{Ex} des $\hand{\xi} \in G(k)$ tels que
$\underline{P}_{\hand{\xi}}$ ait un point\inserted{\unsure{z}} au dessus de x \xout{en lequel on ait bien $\underline{P}$ P soit ramifié}
ne satisfaisant \uline{pas} l'ensemble de conditions suivantes: $\underline{P}_{\hand{\xi}}$ est lisse en
z, \xout{le morphisme $\underline{P}$ P est non ramifié en z, et l'image de l'espace}
l'application tangente à $\underline{P}_{\hand{\xi}} \rightarrow P$ en z : $T_z(\underline{P}\ ) \rightarrow T_x(P)$ est injective
(i.e. $\underline{P}_{\hand{\xi}} \rightarrow P$ non ramifié en z) et son image est ``transversale'' à V
i.e. sa somme avec V est $T_x(P)$. \uline{Alors} E(x,V) (dont nous savons que
c'est la trace d'un ensemble \xout{constructible} \inserted[typed]{\deleted{fermé}}\inserted{\unsure{constructible}} bien déterminé de G sur G(k))
\uline{est de dimension} $\hand{\leqslant}$ N-n\unsure{\over{+}{-}}1 . - Moyennant cette condition, \xout{xx} l'appli\unsure{:}
cation du critère jacobien et un comptage de dimensions montre que
l'ensemble\inserted{fermé E} des poihts\inserted{x} de $\underline{X}$ tels que $\underline{X} \longrightarrow G$ soit non lisse en x, \uline{ou}
$\underline{P} \longrightarrow$ \unsure{\over{G}{G}} non lisse en $\underline{f}(x)$, ou $\underline{P} \longrightarrow P$ ramifié en $\underline{f}(x)$, est de
dimension $\hand{\leqslant}$ n + (N-n\unsure{\over{+}{-}}1) = N-1 , \hand{(}X étant lisse partout de dimension n\hand{)}.

\margin{94}
